\section{Introduction}
Autonomous agents have achieved remarkable advances across diverse domains, including code generation (e.g., Claude Code~\cite{anthropic2025claude}), scientific discovery (e.g., Biomni~\cite{huang2025biomni}, STELLA~\cite{jin2025stellaselfevolvingllmagent}, and SciMaster~\cite{chai2025scimastergeneralpurposescientificai}), and in-depth research (e.g., Tongyi, OpenAI, and Gemini DeepResearch~\cite{tongyidr,openai2025deepresearch,google2025gemini}), showcasing their immense potential to solve complex, open-ended real-world tasks end to end. 

%然而，在特定任务上预训练完的agent不能和环境交互积累经验，持续进化，从而在困难、未训练的任务上表现很差。现有的基于梯度的训练方法存在训练开销大和灾难性遗忘的问题。并且，很多SOTA的LLM模型都是闭源的，让更新参数的梯度优化不可能实现。
However, pretrained agents remain static with frozen parameters, fundamentally precludes learning from on-the-fly experience concluded during trial-and-error, causing substantial performance degradation on challenging or untrained tasks~\cite{wu2023autogen,liu2023agentbench}.
While gradient-based learning has underpinned model optimization~\cite{rumelhart1986learning,ruder2016overview,hu2022lora}, it is inherently ill-suited for continuous agent evolution due to three key obstacles:
(i) back-propagation incurs prohibitive computational costs~\cite{goodfellow2013empirical,cai2025training};
(ii) parameter-centric paradigm suffers from catastrophic forgetting~\cite{abbas2023loss,hadi2023large}, failing to effectively learn form accumulated experiences over times;
(iii) the closed-source nature of most state-of-the-art LLMs~\cite{google2025gemini2_5_pro,openai2025gpt5,anthropic2025sonnet4.5,anthropic2025claude4,xai2025grok4} renders direct parameter optimization infeasible.

%现有self-evolving agent的工作通过环境交互，进化prompt，workflow，tools，提升agent表现。
Recent self-evolving paradigms attempt to overcome this limitation by evolving \textit{non-parametric} components (\textit{i.e.,} prompts~\cite{yuksekgonul2024textgrad,zhang2025agentic}, workflows~\cite{lin2025se,zhang2025multi,shang2024agentsquare}, and tools~\cite{qian2023creator,qiu2025alita}) through environmental feedback and textual gradient~\cite{yuksekgonul2024textgrad,shinn2023reflexion,zhang2024revolve}. 
%然而，只能针对特定任务优化，不具备continuous evolution
%（感觉这里介词换成over会不会比towards更好？ --xinyuan）
%over语法不对，现在我们要说的是阻碍evolution (GET)
Nevertheless, these approaches encounter three vital bottlenecks towards continuous agent evolution: 
%缺一个逻辑，为什么（这里缺的是什么逻辑呢？）
(i) their evolving components are task-specific, thus they cannot generalize across tasks, limiting cross-task evolution.
(ii) Their evolving components (\textit{i.e.} prompts, workflows, and tools) are not continuously scalable, thus they can only leverage a limited set of experiences, preventing performance from scaling with accumulated knowledge.
(iii) Their evolving components are model-specific, thus new agents must interact from scratch during deployment, unable to continuously evolve from agents' previous off-policy experiences, leading to substantial computational overhead.

% Our methods
To address these limitations, we propose \textbf{F}orward \textbf{L}earning from \textbf{Ex}perience~(\textbf{\method{}}), a novel learning paradigm that shifts learning from modifying \textit{model parameters} to constructing and leveraging an evolvable \textit{experience library} (Fig.\ref{fig:main_method_overview}).
%显式对比 + 三个要点 (explore, update, guide future)
In contrast to traditional gradient-based learning paradigms, which optimize through both forward- and backward-propagation, \method{} involves mere forward passes in three stages: 
(i) exploring forwardly to acquire extensive problem-solving experiences;
(ii) evolving the experience library via the updater;
(iii) guiding forward reasoning by utilizing pertinent experiences from the library.

%% More details of this method
%需要层层递进地解释为什么我们没有上面的问题（exp-lib是跨session持久化的，除此之外我们还发现了两大特性）
Hence, \method{} overcomes the prevailing bottlenecks with three profound properties for continuous evolution: 
(i) \method{} evolves a persistent experience library that continuously aggregates cross-task experiences, ensuring that knowledge is retained and available for future use.
(ii) The library continuously expands and refines, enabling agents to progressively acquire deeper insights and knowledge. This accumulation allows agent performance to scale with experience.
(iii) \method{} stores strategies distilled from experiences in a semantic form, making the library both transferable and interpretable. This design enables seamless inheritance across agents, avoiding redundant and costly re-training.
%(i) The evolving object of \method{} is the persistent experience library, which can continuously collect cross-task experience, guaranteeing experience accumulation and preservation for future utilization.
%(ii) The library can continuously expand and refine, enabling agents to progressively acquire deeper insights and knowledge, scaling their cognitive capabilities with accumulated experiences, achieving continuous evolution.
%(iii) \method{} stores strategies inducted from experiences in a semantic form, endowing the experience library with transferability and interpretability, enabling seamless inheritance across agents without redundant and costly re-training.
%(iii) The inheritance of the experience library, where learned experience libraries can be seamlessly inherited across agents. \method{} overcomes this bottleneck by leveraging the great transferability and interpretability of inducted textual signals.

% Experimental Results
To validate the general effectiveness of \method{} paradigm, we conduct extensive experiments across diverse challenging scientific domains, including Olympiad-level mathematics (AIME25~\cite{aime2025problems}), chemical retrosynthesis (USPTO50k~\cite{schneider2016s}), and protein fitness prediction (ProteinGym~\cite{notin2023proteingym}). \method{} demonstrates substantial and consistent improvements on these tasks, from $40\%$ to $63\%$ on AIME25, $20\%$ to $30\%$ on USPTO50k, and $46\%$ to $60\%$ on ProteinGym, exhibiting great enhancement in the capacity of reasoning and knowledge leverage. 

% Contribution summary
In summary, our contributions are as follows:
\begin{itemize}
\item We propose \textbf{\method{}}, a new learning paradigm for the agentic era. It redefines learning as a forward exploration and experience distillation process, enabling LLM agents to continuously evolve through accumulated experience without gradient-based tuning.
\item We provide a comprehensive framework for \method{}, including a unified mathematical formulation with theoretical justifications, a practical instantiation with concrete mechanisms, and an empirical demonstration of its effectiveness on diverse scientific benchmarks.
\item We identify and empirically validate a \textit{scaling law} of the experience library, showing that agent performance scales predictably with accumulated experience. %and revealing a path towards a collaborative experience ecosystem.
We also introduce and demonstrate the principle of \textit{experience inheritance}, where distilled experience can be transferred between agents in a plug-and-play manner, enabling instant knowledge assimilation and bypassing redundant learning.
\end{itemize}